\uuid{394p}
\exo7id{7518}
\titre{exo7 7518}
\auteur{mourougane}
\organisation{exo7}
\datecreate{2021-08-10}
\isIndication{false}
\isCorrection{false}
\chapitre{Nombres complexes}
\sousChapitre{Forme cartésienne, forme polaire}
\cpge{mpsi-mp2i.nombres-complexes.b, pcsi.nombres-complexes.b, ptsi.nombres-complexes.b, mpsi-mp2i.nombres-complexes.a, pcsi.nombres-complexes.a, ptsi.nombres-complexes.a, bcpst1.nombres-complexes.a}
\module{Algèbre}
\niveau{L1}
\difficulte{}

\contenu{
\begin{enumerate}
    \item \question{Soit $z=a+ib$ un nombre complexe non nul.
    Déterminer la partie réelle et la partie imaginaire de l'inverse de $z$.}
    \item \question{Soit $t$ un nombre réel et $z=a+ib$ un nombre complexe.
    Déterminer la partie réelle et la partie imaginaire de $\frac{z-t}{z+t}$.}
    \item \question{Soit $z$ un nombre complexe. Démontrer que $z$ est réel si et seulement si $\overline{z}=z$.}
\end{enumerate}
}
